\documentclass[twocolumn,pra,aps,superscriptaddress,floatfix,10pt]{revtex4-2}

\usepackage{amsmath}
\usepackage{graphicx}
\usepackage{xcolor}
\usepackage{bm}
\usepackage[colorlinks=true,linkcolor=black,citecolor=black,urlcolor=blue]{hyperref}

\newcommand{\etal}{\textit{et~al.}}

\begin{document}

\title{Nonlinear Breit-Rabi amplification of sub-ppm field inhomogeneity in positronium hyperfine spectroscopy}

\author{Ernest Darell Zerme\~no}
\email{0244552@up.edu.mx}
\affiliation{Universidad Panamericana, Guadalajara, M\'exico}

\date{\today}

\begin{abstract}
In their 1975 precision measurement of the positronium ground-state hyperfine structure (HFS), Mills and Bearman characterized the magnetic field gradients across their microwave cavity by nuclear magnetic resonance (NMR) probing and reported a sub-ppm spatial offset $\langle\delta B/B_0\rangle = +0.44$\,ppm averaged over the positronium distribution, which they applied as a linear correction to the central field. We show that propagating this same offset through the full nonlinear Breit-Rabi inversion---a step not performed in the original analysis or any subsequent work---produces an additional bias on the extracted HFS interval of $\delta\Delta_\mathrm{HFS} = -0.22 \pm 0.10$\,MHz, with no free parameters. The underlying mechanism is a sensitivity coefficient $d\Delta_\mathrm{HFS}/dB_\mathrm{res} \approx 4.5 \times 10^5$\,MHz/T near the Breit-Rabi crossing, which converts ppm-level field offsets into MHz-level HFS shifts. The same calculation applied to the Ritter \etal\ (1984) measurement gives $-0.18$ to $-0.30$\,MHz; for the Ishida \etal\ (2014) superconducting solenoid (1.5\,ppm uniformity), $|\delta\Delta_\mathrm{HFS}| < 0.3$\,MHz, consistent with their agreement with QED. The bias preserves the Lorentzian line shape to better than $10^{-4}$ of the peak height. Its pressure dependence is bounded by the unknown variation of the transverse positronium distribution with gas density: in the limiting cases, either the full $-0.22$\,MHz or a larger axial-only component of $-0.73$\,MHz survives the standard zero-density extrapolation. In both limits, a non-negligible bias persists after extrapolation. The Lorentzian shape preservation explains why this systematic escaped detection by line-shape diagnostics for four decades.
\end{abstract}

\maketitle

%% ====================================================================
\section{Introduction}
\label{sec:intro}
%% ====================================================================

Positronium (Ps), the bound state of an electron and a positron, is the lightest purely leptonic atom and provides a uniquely clean system for testing bound-state quantum electrodynamics (QED)~\cite{Adkins2022,Karshenboim2005}. The current theoretical prediction for the ground-state hyperfine structure (HFS) interval, including corrections through $O(m\alpha^6)$ and partial $O(m\alpha^7)$ contributions, is $\Delta_\mathrm{HFS}^\mathrm{QED} = 203\,391.69 \pm 0.41$\,MHz~\cite{Kniehl2000,Melnikov2001,Baker2014}. The two most precise measurements, by Mills \& Bearman~\cite{Mills1975,Mills1983} ($203\,387.5 \pm 1.6$\,MHz) and Ritter \etal~\cite{Ritter1984} ($203\,389.10 \pm 0.74$\,MHz), both used the indirect Zeeman method in gas-filled microwave cavities within conventional electromagnets. Their weighted average of $203\,388.65 \pm 0.67$\,MHz~\cite{Adkins2022} lies $3.9\sigma$ below QED. A more recent measurement by Ishida \etal~\cite{Ishida2014}, using a superconducting solenoid with 1.5\,ppm field uniformity, obtained $203\,394.2 \pm 2.1$\,MHz---consistent with QED but with insufficient precision to resolve the discrepancy.

The role of magnetic field inhomogeneity as a possible systematic in the conventional-magnet measurements has been recognized qualitatively for over a decade. Miyazaki~\cite{Miyazaki2010} identified field non-uniformity and non-thermalization as ``probable common systematic errors.'' Asai \etal~\cite{Asai2012} noted the same concern in designing a direct 203\,GHz measurement. The ETH Zurich group~\cite{Heiss2018} explicitly designed their experiment to be ``free of systematic effects found in earlier experiments, namely the inhomogeneity in static magnetic fields.'' Mills and Bearman themselves measured the gradients across their cavity by NMR and reported a sub-ppm spatial offset, which they applied as a linear correction to the central magnetic field~\cite{Mills1975}. However, the offset was never propagated through the nonlinear Breit-Rabi inversion that converts the resonance field into the extracted HFS interval. Although Mills~\cite{Mills1983} later revisited the 1975 measurement to address line-shape asymmetry from singlet annihilation, that reanalysis treated only the non-Hermitian contribution to the line shape; it did not reconsider the inhomogeneity-induced bias discussed here. We perform that propagation for the first time.

We show in Sec.~\ref{sec:mechanism} that a sensitivity coefficient $d\Delta_\mathrm{HFS}/dB_\mathrm{res} \approx 4.5 \times 10^5$\,MHz/T near the Breit-Rabi crossing amplifies sub-ppm field offsets into MHz-level HFS shifts. Section~\ref{sec:mills} applies this to the Mills measurement using the published gradients and Ps distribution, with no free parameters, obtaining $\delta\Delta_\mathrm{HFS} = -0.22 \pm 0.10$\,MHz. Sections~\ref{sec:ritter} and~\ref{sec:ishida} apply the same calculation to Ritter (parameter-bounded) and Ishida (negative control). Section~\ref{sec:invisible} identifies two general properties---Lorentzian shape preservation and pressure independence---that explain why this systematic escaped detection for four decades.

%% ====================================================================
\section{Mechanism: Breit-Rabi amplification}
\label{sec:mechanism}
%% ====================================================================

\subsection{Positronium in a magnetic field}

The ground-state Hamiltonian of positronium in a magnetic field $B$ along $\hat{z}$ is $H = H_\mathrm{HFS} + H_\mathrm{Z}$, where $H_\mathrm{HFS} = (\Delta/4)\,\bm{\sigma}_e \cdot \bm{\sigma}_p$ is the hyperfine interaction and $H_\mathrm{Z} = g_e \mu_B B\,(S_{ez} - S_{pz})/\hbar$ is the Zeeman interaction. Here $g_e = 2.002\,319\,304$ and $\mu_B/h = 13\,996.245$\,MHz/T. The $m = \pm 1$ triplet states are exactly unaffected by the field; only the $m = 0$ subspace is mixed, giving eigenvalues
\begin{align}
  E_{m=\pm 1} &= \Delta/4, \label{eq:Em1} \\
  E_{\pm} &= -\Delta/4 \pm \tfrac{1}{2}\sqrt{\Delta^2 + 4V^2}, \label{eq:Epm}
\end{align}
where $V = g_e \mu_B B$. In our numerical implementation, we diagonalize the full $4 \times 4$ Hamiltonian matrix, which agrees with Eqs.~(\ref{eq:Em1})--(\ref{eq:Epm}) to better than $10^{-11}$\,MHz. The measured transition frequency is
\begin{equation}
  \nu_\mathrm{trans}(B) = E_+ - E_{m=+1} = \frac{\Delta}{2}\bigl(\sqrt{1 + x^2} - 1\bigr),
  \label{eq:nutrans}
\end{equation}
where $x = 2V/\Delta$. Figure~\ref{fig:zeeman} shows the four eigenvalues as a function of $B$.

\begin{figure}[htbp]
  \centering
  \includegraphics[width=\columnwidth]{fig1_zeeman_eigenvalues.png}
  \caption{Zeeman eigenvalues of the positronium ground state. The $m = \pm 1$ states (green, orange) are exactly degenerate at $\Delta/4$ for all $B$. The purple arrow indicates the measured transition $E_+ - E_{m=+1}$.}
  \label{fig:zeeman}
\end{figure}

\subsection{The amplification factor}
\label{sec:amplification}

In the indirect Zeeman method, one locates the central field $B_\mathrm{res}$ at which $\nu_\mathrm{trans}(B_\mathrm{res}) = \nu_\mathrm{MW}$ for a fixed microwave frequency $\nu_\mathrm{MW}$, and inverts Eq.~(\ref{eq:nutrans}) to obtain $\Delta_\mathrm{HFS}$. If the magnetic field is non-uniform, each positronium atom samples a local field $B_i = B_0(1 + \delta_i)$. To leading order, the averaged signal peaks at an apparent field shifted by $\delta B_\mathrm{res} \approx -\langle\delta\rangle B_0$ relative to the true resonance field.

The effect on the extracted $\Delta_\mathrm{HFS}$ is obtained by implicit differentiation of Eq.~(\ref{eq:nutrans}) at fixed $\nu_\mathrm{MW}$:
\begin{equation}
  \frac{d\Delta_\mathrm{HFS}}{dB_\mathrm{res}} = -\frac{\partial\nu_\mathrm{trans}/\partial B}{\partial\nu_\mathrm{trans}/\partial\Delta}.
  \label{eq:amplification}
\end{equation}
At the operating point of the Mills experiment ($B \approx 0.925$\,T, $\nu_\mathrm{MW} = 3253$\,MHz), we obtain
\begin{equation}
  \frac{d\Delta_\mathrm{HFS}}{dB_\mathrm{res}} \approx +4.5 \times 10^5~\mathrm{MHz/T}.
  \label{eq:amp_value}
\end{equation}
A spatially averaged offset of just $+0.5$\,ppm at $B_0 = 0.925$\,T ($\delta B \approx 4.6 \times 10^{-7}$\,T) therefore produces $|\delta\Delta_\mathrm{HFS}| \approx 0.2$\,MHz---comparable to the statistical uncertainty of the original measurements. This amplification factor of order $5 \times 10^5$ between the fractional field offset and the extracted HFS shift is generic to indirect Zeeman measurements near the Breit-Rabi crossing.

%% ====================================================================
\section{Reanalysis of Mills (1975)}
\label{sec:mills}
%% ====================================================================

\subsection{Experimental configuration}

Mills \& Bearman~\cite{Mills1975} used a thin silver-plated copper cavity (diameter 112.5\,mm, thickness 12.7\,mm) operating in the TM$_{110}$ mode at $\nu_\mathrm{MW} = 3253$\,MHz with $B_0 = 0.925$\,T. NMR probe measurements at the cavity center yielded $\partial^2 B/\partial x^2 = -14$\,ppm/cm$^2$, $\partial^2 B/\partial y^2 = -28$\,ppm/cm$^2$, and $\partial^2 B/\partial z^2 = +44$\,ppm/cm$^2$ ($z$ along the field). Table~\ref{tab:experiments} summarizes the experimental parameters for all three measurements.

\begin{table}[b]
\caption{Experimental parameters for the three HFS measurements.}
\label{tab:experiments}
\begin{ruledtabular}
\begin{tabular}{lccc}
 & Mills~\cite{Mills1975} & Ritter~\cite{Ritter1984} & Ishida~\cite{Ishida2014} \\
\hline
$B_0$ (T) & 0.925 & 0.79 & 0.866 \\
$\nu_\mathrm{MW}$ (MHz) & 3253 & 2384 & 2856.6 \\
Cavity $R$ (mm) & 56.25 & 76.8 & 64 \\
Cavity $L$ (mm) & 12.7 & 38.1 & 100 \\
Magnet type & Electromagnet & Electromagnet & SC solenoid \\
$B$ uniformity & Measured\footnotemark[1] & $\sim 0.5$\,ppm\footnotemark[2] & 1.5\,ppm\footnotemark[3] \\
\end{tabular}
\end{ruledtabular}
\footnotetext[1]{NMR-measured second derivatives reported in Ref.~\cite{Mills1975}.}
\footnotetext[2]{Post-shimming weighted standard deviation~\cite{Ritter1984}.}
\footnotetext[3]{RMS uniformity in a $40 \times 100$\,mm region~\cite{Ishida2014}. An earlier value of 0.9\,ppm was reported in Ref.~\cite{Ishida2011}.}
\end{table}

\subsection{Field model and spatial distribution}

In cylindrical symmetry, the azimuthal average gives $\partial^2 B/\partial r^2 = (\partial^2 B/\partial x^2 + \partial^2 B/\partial y^2)/2 = -21$\,ppm/cm$^2$. For the Gaussian transverse distribution measured by Mills ($\sigma_r = 2.44$\,mm), the azimuthal average is exact: the contribution $\langle\tfrac{1}{2}(\alpha_x x^2 + \alpha_y y^2)\rangle = \tfrac{1}{2}(\alpha_x + \alpha_y)\sigma^2 = \alpha_r \sigma^2$ is independent of the individual values of $\alpha_x$ and $\alpha_y$ when $\langle x^2\rangle = \langle y^2\rangle = \sigma^2$. The field profile is
\begin{equation}
  B(r,z) = B_0\left[1 + \frac{1}{2}\alpha_r r^2 + \frac{1}{2}\alpha_z z^2\right],
  \label{eq:mills_field}
\end{equation}
where $\alpha_r = -0.21$\,m$^{-2}$ and $\alpha_z = +0.44$\,m$^{-2}$. This model has no free parameters---it is fully determined by the published experimental data. We note that the measured gradients satisfy $\nabla^2 B_z = 0$ only approximately: the sum $\partial^2 B/\partial x^2 + \partial^2 B/\partial y^2 + \partial^2 B/\partial z^2 = +2$\,ppm/cm$^2$ has a residual of $2.3\%$ of the total variation. If one imposes the Maxwell constraint $\nabla^2 B_z = 0$ exactly, the axial curvature becomes $+42$\,ppm/cm$^2$ ($\alpha_z = +0.42$\,m$^{-2}$), yielding $\langle\delta B/B_0\rangle = +0.44$\,ppm---in exact agreement with the value reported by Mills \& Bearman~\cite{Mills1975}. Our baseline calculation uses the measured $\alpha_z = +0.44$\,m$^{-2}$; the Maxwell-constrained alternative is included in the systematic uncertainty.

The positronium transverse distribution was measured by slit scans~\cite{Mills1975}: Gaussian with $\sigma_r = 2.44$\,mm. The axial distribution is uniform within the 12.7\,mm cavity thickness, weighted by the detection efficiency function $w(z) = 1 - b|z|$ with $b = 1.4$\,cm$^{-1}$ (Ref.~\cite{Mills1975}, Eq.~4), which represents the position-dependent coincidence detection efficiency for 2$\gamma$ annihilation events. The weighting remains positive throughout the cavity ($b L/2 = 0.89 < 1$).

\subsection{Simulation procedure}

We numerically reproduce the complete experimental procedure:
\begin{enumerate}
  \item Sample $N = 5 \times 10^5$ positronium positions $(r_i, z_i)$ from the Gaussian transverse distribution and uniform axial distribution.
  \item For each scan field $B_0$, compute the local field $B_i = (B_0/B_0^\mathrm{nom}) \times B(r_i, z_i)$.
  \item Compute the local transition frequency using a pre-computed cubic interpolation table (20\,000 points).
  \item The signal at each scan point is the detection-weighted Lorentzian average:
  \begin{equation}
    S(B_0) = \frac{\sum_i w(z_i)\, L_i(B_0)}{\sum_i w(z_i)},
    \label{eq:signal}
  \end{equation}
  where $L_i = (\gamma/2)^2/[(\nu_{\mathrm{trans},i} - \nu_\mathrm{MW})^2 + (\gamma/2)^2]$ and $\gamma = 5$\,MHz.
  \item Fit $S(B_0)$ to a Lorentzian in $B$ to extract $B_\mathrm{res}$.
  \item Extract $\Delta_\mathrm{HFS}$ by numerically inverting Eq.~(\ref{eq:nutrans}).
\end{enumerate}

A baseline computed with the identical pipeline using a uniform field ($B = B_0$ everywhere) is subtracted to isolate the inhomogeneity effect from the nonlinear fitting bias ($\approx -0.03$\,MHz):
\begin{equation}
  \delta\Delta_\mathrm{HFS} = \Delta_\mathrm{HFS}^\mathrm{extracted} - \Delta_\mathrm{HFS}^\mathrm{uniform\;control}.
  \label{eq:shift}
\end{equation}
This baseline subtraction is robust: the corrected result is stable to $< 0.001$\,MHz for scan windows ranging from $\pm 3\gamma$ to $\pm 16\gamma$ and from 50 to 800 scan points.

\subsection{Result}

Using the published gradients and the measured positronium distribution with detection weighting, we obtain
\begin{equation}
  \delta\Delta_\mathrm{HFS}^\mathrm{Mills} = -0.22 \pm 0.10~\mathrm{MHz},
  \label{eq:mills_result}
\end{equation}
where the uncertainty is the quadrature sum of: (i) Monte Carlo statistical uncertainty, $\pm 0.001$\,MHz (from 20 independent random seeds); (ii) NMR gradient measurement uncertainty, $\pm 0.02$\,MHz (the difference between using the measured $\alpha_z = +0.44$\,m$^{-2}$ and the Maxwell-constrained $\alpha_z = +0.42$\,m$^{-2}$, which changes $\langle\delta B/B_0\rangle$ from $+0.52$ to $+0.44$\,ppm and $\delta\Delta$ from $-0.22$ to $-0.18$\,MHz); (iii) the detection weight parameter $b$, $\pm 0.09$\,MHz (for a $\pm 10\%$ variation around the published $b = 1.4$\,cm$^{-1}$; the sensitivity is steep, producing the quoted $\pm 0.09$\,MHz variation); (iv) TM$_{110}$ microwave mode weighting, $\pm 0.014$\,MHz (verified numerically; see Sec.~\ref{sec:limitations}); and (v) pressure dependence of the radial distribution, $\pm 0.02$\,MHz (at the operating conditions). The dominant uncertainty is the detection weight $b$; if $b$ were known to $\pm 5\%$, the total uncertainty would reduce to $\pm 0.06$\,MHz. The mean field offset over the positronium distribution is $\langle\delta B/B_0\rangle = +0.52$\,ppm using the measured gradients, or $+0.44$\,ppm if the Maxwell constraint is imposed---reproducing exactly the value reported by Mills \& Bearman~\cite{Mills1975}.

The result scales linearly with the gradient magnitude to better than $0.1\%$ (tested up to $4\times$ the measured gradients), confirming that the first-order mechanism described in Sec.~\ref{sec:amplification} dominates. The analytical estimate from the amplification factor, $\delta\Delta_\mathrm{HFS} \approx (d\Delta/dB_\mathrm{res}) \times \langle\delta B\rangle \approx 4.5 \times 10^5 \times (-0.52 \times 10^{-6} \times 0.925) = -0.22$\,MHz, agrees with the full numerical result to $1\%$.

Figure~\ref{fig:resonance} shows the simulated resonance curve.

\begin{figure}[htbp]
  \centering
  \includegraphics[width=\columnwidth]{resonance_curves_corrected.png}
  \caption{Upper panel: simulated Zeeman resonance for the Mills configuration with NMR-measured gradients (solid blue) and Lorentzian fit (dashed red), compared to the uniform-field control (green). The annotated $\delta\Delta_\mathrm{HFS}$ is baseline-subtracted. Lower panel: fit residuals normalized to the peak height, confirming $< 10^{-4}$ deviation from a perfect Lorentzian.}
  \label{fig:resonance}
\end{figure}

%% ====================================================================
\section{Reanalysis of Ritter (1984)}
\label{sec:ritter}
%% ====================================================================

Ritter \etal~\cite{Ritter1984} used a Varian V4012A electromagnet (pole diameter 304.8\,mm, gap 66.7\,mm)~\cite{Theriot1970} with careful shimming. The field is parameterized as
\begin{equation}
  B(r,z) = B_0\left[1 + c_2\,\frac{2z^2 - r^2}{R_\mathrm{ref}^2}\right],
  \label{eq:ritter_field}
\end{equation}
with $R_\mathrm{ref} = 33.35$\,mm (half the pole gap). Ritter \etal\ reported a weighted standard deviation of 0.5\,ppm after shimming, corresponding to $c_2 \approx 3$--$5$\,ppm. We use the actual cavity dimensions ($R_\mathrm{cav} = 76.8$\,mm, $L_\mathrm{cav} = 38.1$\,mm) and positronium confinement region ($R \approx 12.7$\,mm, uniform). Table~\ref{tab:results} presents the results.

\begin{table}[t]
\caption{Systematic shift $\delta\Delta_\mathrm{HFS}$ for each experimental configuration.}
\label{tab:results}
\begin{ruledtabular}
\begin{tabular}{llr}
Configuration & Parameters & $\delta\Delta_\mathrm{HFS}$ \\
 & & (MHz) \\
\hline
Mills (meas.\ gradients) & Published data & $-0.22 \pm 0.10$ \\
Ritter ($c_2 = 3$\,ppm) & Post-shimming & $-0.18$ \\
Ritter ($c_2 = 5$\,ppm) & Post-shimming & $-0.30$ \\
Ritter ($c_2 = 10$\,ppm) & Pre-shimming & $-0.60$ \\
Ritter ($c_2 = 20$\,ppm) & Pre-shimming & $-1.19$ \\
Ishida SC (1.5\,ppm) & Uniform Ps & $+0.30$ \\
Ishida SC (1.5\,ppm) & $\lambda = 30$\,mm & $-0.09$ \\
\end{tabular}
\end{ruledtabular}
\end{table}

For the post-shimming range $c_2 = 3$--$5$\,ppm, the bias is $-0.18$ to $-0.30$\,MHz. The shift scales linearly with $c_2$, as expected from the first-order mechanism. Figure~\ref{fig:ritter_scan} shows the systematic shift as a function of $c_2$.

\begin{figure}[htbp]
  \centering
  \includegraphics[width=\columnwidth]{ritter_c2_scan.png}
  \caption{Systematic shift for the Ritter configuration with actual cavity dimensions and positronium confinement. Green shading: post-shimming range ($c_2 = 2$--$5$\,ppm); orange shading: pre-shimming range. The dashed line marks the Ritter discrepancy with QED ($-2.59$\,MHz) for reference.}
  \label{fig:ritter_scan}
\end{figure}

%% ====================================================================
\section{Verification with Ishida (2014)}
\label{sec:ishida}
%% ====================================================================

The Ishida experiment~\cite{Ishida2014} used a superconducting solenoid with 1.5\,ppm RMS residual non-uniformity in a $40 \times 100$\,mm region, operating at $B_0 = 0.866$\,T and $\nu_\mathrm{MW} = 2856.6$\,MHz. Our simulation gives $\delta\Delta_\mathrm{HFS} = +0.30$\,MHz for uniform Ps and $-0.09$\,MHz for non-thermalized Ps ($\lambda = 30$\,mm). Both values are well below the experimental uncertainty of $\pm 2.1$\,MHz, confirming that the superconducting magnet reduces this systematic to a negligible level.

%% ====================================================================
\section{Two properties that hid the systematic}
\label{sec:invisible}
%% ====================================================================

\subsection{Lorentzian shape preservation}
\label{sec:shape}

The averaged resonance curve $S(B_0)$ [Eq.~(\ref{eq:signal})] retains its Lorentzian shape to remarkable precision. For the Mills configuration, the maximum absolute residual between the inhomogeneous signal and the Lorentzian fit is $9.6 \times 10^{-5}$ of the peak height, with an RMS residual of $5.3 \times 10^{-5}$. This is indistinguishable from the numerical noise of the uniform-field baseline ($9.3 \times 10^{-5}$ and $5.3 \times 10^{-5}$, respectively).

No diagnostic based on the chi-squared of the fit, the presence of asymmetric shoulders, or any other line-shape analysis could have detected this systematic at the precision of the 1975--1984 experiments.

\subsection{Partial pressure independence}
\label{sec:pressure}

The bias can be decomposed into a radial contribution from $\alpha_r$ and an axial contribution from $\alpha_z$:
\begin{equation}
  \langle\delta B/B_0\rangle = \underbrace{\alpha_r \sigma_r^2}_{\text{radial: }-1.25\,\text{ppm}} + \underbrace{\tfrac{1}{2}\alpha_z \langle z^2\rangle_w}_{\text{axial: }+1.77\,\text{ppm}}.
  \label{eq:decomposition}
\end{equation}
The axial term depends only on the detection geometry and cavity thickness---not on the gas pressure. The radial term, however, depends on $\sigma_r$, which varies with pressure through the positron stopping range and subsequent diffusion. Because the two terms have opposite signs and partially cancel at the measured $\sigma_r = 2.44$\,mm, even moderate pressure dependence of $\sigma_r$ can change the net bias significantly: for $\sigma_r$ varying over 1.5--5.0\,mm, the total $\delta\Delta_\mathrm{HFS}$ ranges from approximately $-0.5$ to $+1.4$\,MHz.

The standard zero-density extrapolation removes the pressure-\emph{dependent} part of the bias. If $\sigma_r$ is approximately constant across the pressures used in the extrapolation, the full $-0.22$\,MHz bias survives intact. If $\sigma_r$ varies with gas density $n$, the effect on the extrapolation depends on the functional form of $\sigma_r(n)$ and can be substantial: since the two terms partially cancel at the nominal $\sigma_r$, even moderate variation can change the magnitude or sign of the bias at the extrapolated $n \to 0$ intercept. The axial component alone ($\delta\Delta \approx -0.73$\,MHz) is strictly pressure-independent and survives any extrapolation procedure. Characterizing $\sigma_r$ as a function of gas density, which was not reported in the original experiment, is necessary to determine the fraction of the total bias that survives the extrapolation.

\subsection{Generality}

The Lorentzian shape preservation, together with the partial pressure independence of the axial component, characterizes a class of systematics that largely escape both line-shape diagnostics and pressure-extrapolation controls. Any inhomogeneity-induced shift that is linear in the mean spatial offset of the field and does not couple to gas-density-dependent dynamics falls into this class. The amplification by the nonlinear Breit-Rabi inversion, which is generic to indirect Zeeman methods, converts what appears to be a ``sub-ppm'' effect into a MHz-level bias on the extracted splitting.

%% ====================================================================
\section{Limitations and the joint thermalization problem}
\label{sec:limitations}
%% ====================================================================

Our calculation isolates a single systematic---the field-inhomogeneity bias propagated through the Breit-Rabi inversion---under the simplifying assumption of a thermalized positronium spatial distribution. In reality, the same atoms that sample the non-uniform field are also subject to incomplete thermalization, which Ishida \etal~\cite{Ishida2014} estimated as a separate $\sim 2$\,MHz systematic in their own apparatus.

These two corrections are not independent. The non-thermalized positronium population samples the magnetic field with a different spatial weight than the thermalized population, and the velocity distribution that drives the gas-density shift depends on the local field through the position-dependent formation dynamics. A correct joint correction requires a self-consistent simulation of thermalization in a non-uniform field, which is beyond the scope of the present work.

We do not combine our field-inhomogeneity result with the published Ishida thermalization estimate to produce a single corrected significance for the historical discrepancy. Such a combination would require: (i) re-deriving the thermalization correction under the gas conditions of Mills (isobutane) and Ritter (N$_2$) rather than transferring Ishida's number, (ii) computing the joint distribution of position and velocity, and (iii) propagating the correlated uncertainty.

We have verified the following additional effects numerically. (i) The TM$_{110}$ microwave mode of the Mills cavity has a spatially varying field amplitude $|B_\mathrm{MW}(\mathbf{r})|^2$ that weights the transition rate. Including this weight in the simulation changes the result by $-0.014$\,MHz ($6.6\%$), consistent with the fact that the positronium distribution ($\sigma_r = 2.44$\,mm) occupies less than $3\%$ of the mode's radial scale ($R = 56.25$\,mm), where the mode is nearly uniform. (ii) The non-Hermitian corrections from singlet annihilation, which Mills~\cite{Mills1983} showed produce a line-center shift of $\sim 0.5$\,MHz in $\nu_\mathrm{trans}$, produce a much larger shift of approximately $75$\,MHz in the extracted $\Delta_\mathrm{HFS}$. This amplification occurs because the $0.5$\,MHz shift in transition frequency is magnified by the factor $|d\Delta/d\nu| = 1/|\partial\nu/\partial\Delta| \approx 65$ inherent in the Breit-Rabi inversion, with the remaining factor of $\sim 2$ arising from the non-Hermitian asymmetry of the mixed state near the Breit-Rabi crossing, where the singlet admixture (and thus the annihilation-induced frequency pull) varies quadratically with field. Crucially, this shift is identical for uniform and non-uniform fields; under baseline subtraction, the residual is $< 0.001$\,MHz, confirming complete cancellation. (iii) The difference between our $\langle\delta B/B_0\rangle = +0.52$\,ppm and the $+0.44$\,ppm of Mills \& Bearman is almost entirely explained by the $+2$\,ppm/cm$^2$ Laplacian residual in the NMR measurements: imposing $\nabla^2 B_z = 0$ exactly ($\alpha_z = +0.42$\,m$^{-2}$ instead of $+0.44$) recovers $+0.44$\,ppm. This NMR measurement uncertainty is included in the error budget as $\pm 0.02$\,MHz. (iv) The pressure dependence of $\sigma_r$, which was not characterized in the original experiment, constitutes an irreducible uncertainty in the fraction of the bias that survives the zero-density extrapolation (see Sec.~\ref{sec:pressure}).

The claims of this work are made at three distinct levels of robustness:
\begin{enumerate}
  \item \textbf{Parameter-free.} The field-inhomogeneity contribution to the Mills HFS value is $-0.22 \pm 0.10$\,MHz using only the published gradients and detection weighting.
  \item \textbf{Parameter-bounded.} For the Ritter measurement, the contribution lies in the range $-0.18$ to $-0.30$\,MHz, bounded by the reported post-shimming uniformity.
  \item \textbf{Mechanism-level.} The Lorentzian shape preservation is a general property of the Breit-Rabi amplification mechanism. The fraction of the bias surviving the zero-density extrapolation depends on the unknown $\sigma_r(n)$; the pressure-independent axial component alone contributes $-0.73$\,MHz.
\end{enumerate}

%% ====================================================================
\section{Implications for current and future measurements}
\label{sec:future}
%% ====================================================================

Our results imply that any future precision measurement of the ground-state positronium HFS using the indirect Zeeman method must verify the magnetic field uniformity at the $\sim 1$\,ppm level and demonstrate complete Ps thermalization. The importance of cavity-related systematics has been underscored by recent work at UCL: Akopyan \etal~\cite{Akopyan2021} showed through numerical simulations that microwave reflections produce frequency-dependent power distributions that shift extracted frequencies by several MHz, and Sheldon \etal~\cite{Sheldon2023} confirmed this by redesigning the microwave environment, reducing a $4.5\sigma$ fine-structure discrepancy~\cite{Gurung2020} to $1.3\sigma$. The ongoing programs at ETH Zurich~\cite{ThorpeWoods2026} and the Ramsey spectroscopy techniques at UCL~\cite{Newson2026} introduce new sensitivity to field gradients between interaction regions, making field characterization even more critical.

The self-consistent calculation combining field inhomogeneity with positronium thermalization dynamics---accounting for the correlated sampling of non-uniform fields by non-thermalized positronium---is the key remaining step toward determining the joint correction and its impact on the historical discrepancy.

%% ====================================================================
\section{Conclusions}
\label{sec:conclusion}
%% ====================================================================

\begin{enumerate}
  \item[(i)] The Breit-Rabi inversion of an indirect Zeeman HFS measurement amplifies fractional field offsets by a factor $\sim 5 \times 10^5$ near the operating point of the Mills, Ritter, and Ishida experiments.

  \item[(ii)] Propagating the field offset published by Mills and Bearman through this amplification, with no free parameters, gives a contribution of $-0.22 \pm 0.10$\,MHz to the Mills HFS value. The same calculation gives $-0.18$ to $-0.30$\,MHz for the Ritter measurement and $< 0.3$\,MHz for the Ishida superconducting magnet, qualitatively consistent with the observed pattern of agreement with QED.

  \item[(iii)] The bias preserves the Lorentzian fit shape to better than $10^{-4}$ of the peak height, making it invisible to line-shape diagnostics. The non-Hermitian line-center shift from singlet annihilation ($\sim 0.5$\,MHz in $\nu_\mathrm{trans}$, amplified to $\sim 75$\,MHz in the extracted $\Delta_\mathrm{HFS}$) cancels to $< 0.001$\,MHz under baseline subtraction. The fraction of the bias surviving the zero-density extrapolation depends on the unknown pressure dependence of $\sigma_r$; the axial component ($-0.73$\,MHz) is strictly pressure-independent. The TM$_{110}$ microwave mode weight contributes a $6.6\%$ correction ($-0.014$\,MHz).

  \item[(iv)] Applied to the individual measurements, the correction accounts for $\sim 5\%$ of the Mills discrepancy ($-4.2$\,MHz) and $\sim 7$--$12\%$ of the Ritter discrepancy ($-2.6$\,MHz) with QED. Since Mills contributes only $\sim 17\%$ of the statistical weight in the Mills--Ritter average, the net impact on the combined $3.9\sigma$ discrepancy is smaller. A self-consistent simulation combining field inhomogeneity with positronium thermalization dynamics is required to determine the joint correction.
\end{enumerate}

%% ====================================================================
\begin{acknowledgments}
The simulation code is publicly available at \url{https://github.com/Jefemaestro33/ps-hfs-field-inhomogeneity}.
\end{acknowledgments}
%% ====================================================================

\bibliography{02_paper_refs}

\end{document}
